%%%%%%%%%%%%%%%%%%%%%%%%%%%%%%%%%%%%%%%%
% Compact Laboratory Book
% LaTeX Template
% Version 1.0 (4/6/12)
%
% This template has been downloaded from:
% http://www.LaTeXTemplates.com
%
% Original author:
% Joan Queralt Gil (http://phobos.xtec.cat/jqueralt) using the labbook class by
% Frank Kuster (http://www.ctan.org/tex-archive/macros/latex/contrib/labbook/)
%
% License:
% CC BY-NC-SA 3.0 (http://creativecommons.org/licenses/by-nc-sa/3.0/)
%
% Important note:
% This template requires the labbook.cls file to be in the same directory as the
% .tex file. The labbook.cls file provides the necessary structure to create the
% lab book.
%
% The \lipsum[#] commands throughout this template generate dummy text
% to fill the template out. These commands should all be removed when 
% writing lab book content.
%
% HOW TO USE THIS TEMPLATE 
% Each day in the lab consists of three main things:
%
% 1. LABDAY: The first thing to put is the \labday{} command with a date in 
% curly brackets, this will make a new section showing that you are working
% on a new day.
%
% 2. EXPERIMENT/SUBEXPERIMENT: Next you need to specify what 
% experiment(s) and subexperiment(s) you are working on with a 
% \experiment{} and \subexperiment{} commands with the experiment 
% shorthand in the curly brackets. The experiment shorthand is defined in the 
% 'DEFINITION OF EXPERIMENTS' section below, this means you can 
% say \experiment{pcr} and the actual text written to the PDF will be what 
% you set the 'pcr' experiment to be. If the experiment is a one off, you can 
% just write it in the bracket without creating a shorthand. Note: if you don't 
% want to have an experiment, just leave this out and it won't be printed.
%
% 3. CONTENT: Following the experiment is the content, i.e. what progress 
% you made on the experiment that day.
%
%%%%%%%%%%%%%%%%%%%%%%%%%%%%%%%%%%%%%%%%%

%----------------------------------------------------------------------------------------
%	PACKAGES AND OTHER DOCUMENT CONFIGURATIONS
%----------------------------------------------------------------------------------------                               
%\UseRawInputEncoding
\documentclass[fontsize=11pt, % Document font size
                             paper=letter, % Document paper type
                             twoside, % Shifts odd pages to the left for easier reading when printed, can be changed to oneside
                             captions=tableheading,
                             index=totoc,
                             hyperref]{labbook}

%\documentclass[idxtotoc,hyperref,openany]{labbook} % 'openany' here removes the
   
\usepackage[bottom=10em]{geometry} % Reduces the whitespace at the bottom of the page so more text can fit

\usepackage[english]{babel} % English language
\usepackage{lipsum} % Used for inserting dummy 'Lorem ipsum' text into the template

\usepackage[utf8]{inputenc} % Uses the utf8 input encoding
\usepackage[T1]{fontenc} % Use 8-bit encoding that has 256 glyphs

\usepackage[osf]{mathpazo} % Palatino as the main font
\linespread{1.05}\selectfont % Palatino needs some extra spacing, here 5% extra
\usepackage[scaled=.88]{beramono} % Bera-Monospace
\usepackage[scaled=.86]{berasans} % Bera Sans-Serif

\usepackage{booktabs,array} % Packages for tables

\usepackage{amsmath} % For typesetting math
\usepackage{graphicx} % Required for including images
\usepackage{etoolbox}
\usepackage[norule]{footmisc} % Removes the horizontal rule from footnotes
\usepackage{lastpage} % Counts the number of pages of the document
\usepackage{float}
\usepackage[dvipsnames]{xcolor}  % Allows the definition of hex colors
\definecolor{titleblue}{rgb}{0.16,0.24,0.64} % Custom color for the title on the title page
\definecolor{linkcolor}{rgb}{0,0,0.42} % Custom color for links - dark blue at the moment
\usepackage{pgfgantt}
\addtokomafont{title}{\Huge\color{titleblue}} % Titles in custom blue color
\addtokomafont{chapter}{\color{OliveGreen}} % Lab dates in olive green
\addtokomafont{section}{\color{Sepia}} % Sections in sepia
\addtokomafont{pagehead}{\normalfont\sffamily\color{gray}} % Header text in gray and sans serif
\addtokomafont{caption}{\footnotesize\itshape} % Small italic font size for captions
\addtokomafont{captionlabel}{\upshape\bfseries} % Bold for caption labels
\addtokomafont{descriptionlabel}{\rmfamily}
\setcapwidth[r]{10cm} % Right align caption text
\setkomafont{footnote}{\sffamily} % Footnotes in sans serif

\deffootnote[4cm]{4cm}{1em}{\textsuperscript{\thefootnotemark}} % Indent footnotes to line up with text

\DeclareFixedFont{\textcap}{T1}{phv}{bx}{n}{1.5cm} % Font for main title: Helvetica 1.5 cm
\DeclareFixedFont{\textaut}{T1}{phv}{bx}{n}{0.8cm} % Font for author name: Helvetica 0.8 cm

\usepackage[nouppercase,headsepline]{scrpage2} % Provides headers and footers configuration
\pagestyle{scrheadings} % Print the headers and footers on all pages
\clearscrheadfoot % Clean old definitions if they exist

\automark[chapter]{chapter}
\ohead{\headmark} % Prints outer header

\setlength{\headheight}{25pt} % Makes the header take up a bit of extra space for aesthetics
\setheadsepline{.4pt} % Creates a thin rule under the header
\addtokomafont{headsepline}{\color{lightgray}} % Colors the rule under the header light gray

\ofoot[\normalfont\normalcolor{\thepage\ |\  \pageref{LastPage}}]{\normalfont\normalcolor{\thepage\ |\  \pageref{LastPage}}} % Creates an outer footer of: "current page | total pages"

% These lines make it so each new lab day directly follows the previous one i.e. does not start on a new page - comment them out to separate lab days on new pages
\makeatletter
\patchcmd{\addchap}{\if@openright\cleardoublepage\else\clearpage\fi}{\par}{}{}
\makeatother
\renewcommand*{\chapterpagestyle}{scrheadings}

% These lines make it so every figure and equation in the document is numbered consecutively rather than restarting at 1 for each lab day - comment them out to remove this behavior
\usepackage{chngcntr}
\counterwithout{figure}{labday}
\counterwithout{equation}{labday}

% Hyperlink configuration
\usepackage[
    pdfauthor={}, % Your name for the author field in the PDF
    pdftitle={Laboratory Journal}, % PDF title
    pdfsubject={}, % PDF subject
    bookmarksopen=true,
    linktocpage=true,
    urlcolor=linkcolor, % Color of URLs
    citecolor=linkcolor, % Color of citations
    linkcolor=linkcolor, % Color of links to other pages/figures
    backref=page,
    pdfpagelabels=true,
    plainpages=false,
    colorlinks=true, % Turn off all coloring by changing this to false
    bookmarks=true,
    pdfview=FitB]{hyperref}

\usepackage[stretch=10]{microtype} % Slightly tweak font spacing for aesthetics

%\setlength\parindent{0pt} % Uncomment to remove all indentation from paragraphs

%----------------------------------------------------------------------------------------
%	DEFINITION OF EXPERIMENTS
%----------------------------------------------------------------------------------------

% Template: \newexperiment{<abbrev>}[<short form>]{<long form>}
% <abbrev> is the reference to use later in the .tex file in \experiment{}, the <short form> is only used in the table of contents and running title - it is optional, <long form> is what is printed in the lab book itself

\newexperiment{example}[Example experiment]{This is an example experiment}
\newexperiment{example2}[Example experiment 2]{This is another example experiment}
\newexperiment{example3}[Example experiment 3]{This is yet another example experiment}

\newsubexperiment{subexp_example}[Example sub-experiment]{This is an example sub-experiment}
\newsubexperiment{subexp_example2}[Example sub-experiment 2]{This is another example sub-experiment}
\newsubexperiment{subexp_example3}[Example sub-experiment 3]{This is yet another example sub-experiment}

%----------------------------------------------------------------------------------------
\newcommand{\HRule}{\rule{\linewidth}{0.5mm}} % Command to make the lines in the title page

\setlength\parindent{0pt} % Removes all indentation from paragraphs

\begin{document}

%----------------------------------------------------------------------------------------
%	TITLE PAGE
%----------------------------------------------------------------------------------------
%\frontmatter % Use Roman numerals for page numbers

%\begin{center}

%

\title{
\begin{center}
\href{http://www.bradley.edu}{\includegraphics[height=0.5in]{figs/logoBU1-Print}}
\vskip10pt
\HRule \\[0.4cm]
{\Huge \bfseries Meeting Minutes \\[0.5cm] \Large BEMOSS and Its Enhanced Applications}\\[0.4cm] % Degree
\HRule \\[1.5cm]
\end{center}
}
\author{\Huge Brian Lauer \\ \\ \LARGE blauer@mail.bradley.edu \\[2cm]} % Your name and email address
\date{Beginning April 2019} % Beginning date
\maketitle

%\maketitle % Title page

\printindex
\tableofcontents % Table of contents
\newpage % Start lab look on a new page

\begin{addmargin}[0cm]{0cm} % Makes the text width much shorter for a compact look

\pagestyle{scrheadings} % Begin using headers

%----------------------------------------------------------------------------------------
%	LAB BOOK CONTENTS
%----------------------------------------------------------------------------------------
\labday{Friday, April 26, 2019}
%----------------------------------------------------------------------------------------
Dr. Miah will write a research grant proposal under BEMOSS and is due May 3.
\bigbreak\noindent
All documentation will be maintained in a github repository for this project.
I need to create a Github account and send my username to Dr. Miah. A Google Drive with the same name as the Github repository will be created name DocsBEMOSS.
\bigbreak\noindent
I need to talk to Mr. Mattus about getting a laptop for research. It must be installed with Ubuntu administrative privileges. 
\bigbreak\noindent
Send him an email:
\bigbreak\noindent
Hello Mr. Mattus, I will be working on a project with Dr. Miah. For that we would like to request a laptop available in the department if possible so that we can install Ubuntu operating system. 
\bigbreak\noindent
Try to go see him in the afternoon.
\bigbreak\noindent
I need to learn the following Github Bash terminal commands:
\begin{verbatim}
Git add
Git commit -m "message"
Git remote add origin 'url'
\end{verbatim}
To get a better undestanding of what is going on I need to go through Github tutorials.

%\item Upload latex file to Drive
\labday{Thursday, May 02, 2019}
I met with Reece Bachman, Jordan Ingram, and Robert O'Malley in the lab with the intention of filming a video of the BEMOSS installation from start to finish. However, I misinterpreted Dr. Miah's email and thought we were going to install on the current laptop they have in the lab. Instead, I was expected to have received one from Mr. Mattus with Ubuntu installed. Robert showed me the motor setup with the Zigbee module, L298N dual H bridge module, and buck-boost converter. Later, Jordan demonstrated his work with the HVAC controller.

\labday{Friday, May 03, 2019}
I still am in need of a laptop, so I must email Mr. Mattus once again on Monday, May 6.
\bigbreak\noindent
Dr. Miah showed me some of the Github commands to upload tex and pdf files to the Github repo. We added the meeting minutes files \texttt{meeting.tex} and \texttt{meeting.pdf}.

\labday{Friday, May 10, 2019}
Today, I met with Bob and Jordan to film the installation of BEMOSS on my borrowed laptop. After Bob attempted to install BEMOSS twice, he came to the conclusion that I installed the wrong version of Ubuntu on my machine. BEMOSS requires 16.04.5 not 16.04.6 that I had installed before the first installation. After another attempt at installing the software on Ubuntu 16.04.5, the same issue persisted involving three python modules not installing.

\labday{Wednesday, May 22, 2019}
In this meeting with Dr. Miah, he gave me a flash drive with videos of the demonstration of the previous group's setup. Since I and the former students did not end up meeting in the lab to film the video on my borrowed laptop, they used their laptop instead. After unzipping the .tar.gz files, I was unable to view the videos due to some corruption issues. Dr. Miah said he had issues as well. 
\bigbreak\noindent
We discussed when I should work in the lab. He recommended 9 am to 5 pm Monday through Friday which I decided as my work time; however, the hours can be flexible. Likely, I will need to put in some hours outside of this time frame to get fully up to speed on the software.  

\labday{Tuesday, May 28, 2019}
I made some mistakes committing my changes to both the meeting and lab notebooks. My commits were not visible after push them to the repository. I will have to research this further but for now I must use the commands:
\begin{verbatim}
git add .
git commit -m "insert message here"
git pull origin master
git push --force origin master
\end{verbatim}

Look at bemoss presentation (Google BEMOSS slides first link). When I make presentation, no paragraphs on slides. No formal dress is required. 
\medbreak\noindent
For BEMOSS applications, I need to read the paper [cite here Pipattanasomporn2015.pdf]. These applications are listed in the document.

\labday{Friday, May 31, 2019}
We (RAM group) need to setup a common time to be in the lab. Also, after giving my presentation I found that I need to increase my knowledge of BEMOSS greatly and my presentation skills.

\labday{Tuesday, June 4, 2019}
Add references to images. Next Presentation: how to install bemoss. Follow steps in bemoss wiki. Add own thoughts to slides not just copy from wiki. \texttt{201931-ProgressPres.tex}. Next presentation will be 21st of June. 
\bigbreak\noindent
Move all materials from lab into Job 340.
\bigbreak\noindent
BEMOSS work to be done: 
\begin{enumerate}
\item Create functional icon for dc motor in BEMOSS UI to allow toggling on and off
\item Implement simple speed control algorithm within BEMOSS
\item Interface BLDC motor with BEMOSS
\item Researching agent based architecture with BEMOSS
\item Implementing basic machine learning algorithm within BEMOSS
\end{enumerate}
Put personal contributions into list later. 
\bigbreak\noindent
A conference paper will be written in two months with DC motor interface icon in BEMOSS. 
\bigbreak\noindent
Write questions in notebook for meetings. Need to document everything well. 
\bigbreak\noindent
Need to be very confident in BEMOSS. May be presenting papers in conferences. 
\bigbreak\noindent
Improve speaking skills (i.e. talk louder, better eye contact, use figures on slides/poster to better point out what is going on). 
\bigbreak\noindent
Look at Markdown cheat sheet and write more in README in BEMOSS documentation.

\labday{Thursday, June 13, 2019}
Questions for the meeting:
\begin{itemize}
\item What days should I be in the lab?
\item Can I get a deadline for each project? I need to know what I am doing everyday to be productive and not waste anytime.
\item Could I spend some time today building a schedule (at least 2 hours)?
\item Will I have a proper setup for the motor on the third floor?
\end{itemize}

Send an email to Afrash cc'ed to Dr. Miah that I was able to fix the issue with BEMOSS installation after get it working in the lab. We will have Skype meeting with Ashraf discussing further things to add to BEMOSS along with updating BEMOSS with the latest version of Ubuntu. 
\bigbreak\noindent
I will need to update Reece's motor setup. Need to create interface to list all devices available on the network. Need to get previous group's motor setup working. Will be given a new RPi and motor to work on the setup from scratch.
\bigbreak\noindent
Install Skype and send username to Dr. Miah.
\bigbreak\noindent
Need to start writing conference paper, Dr. Miah will share template.
\bigbreak\noindent
May be able to get paid by working with DC motor thus I need to start writing down my hours. May need to setup Bradley account to get paid. 
\bigbreak\noindent
Help Dr. Miah at 4:00 pm with packing items in office. 
\bigbreak\noindent
Do need to by in the lab everyday of the week.
\bigbreak\noindent
Need to have first part (DC motor integration and GUI built) completed by June 28.
\bigbreak\noindent
Under implementation folder, have alternativeInstallation files directory with apache-cassandra database. 
\bigbreak\noindent
Remove files from progressPresentation directory from Github repo and add tex files to parent directory and pdf files to progressPresentation directory.
\bigbreak\noindent
Make sure not to download malware when downloading skype (i.e. don't click on ads).

\labday{Friday, June 21, 2019}
Caleb's presentation:
\bigbreak\noindent
\begin{itemize}
\item Beaglebone Blue - Singleboard computer
\item What is RSS?
\item Recieved Signal Strength needed to track robot's position
\end{itemize}
Might need to create IPE diagram to understand database.  

Amr's Presentation:
\bigbreak\noindent
\begin{itemize}
	\item What is MATLAB used for?
	Controls the robot itself.
	\item What version of Ubuntu are you using on your laptop?
	Dual boot windows, Ubuntu
	\item It is not possible to install ROS in windows. Must use Unix
\end{itemize}
For next presentation, I need to reference my previous presentation and understand the Cassandra database.
\bigbreak\noindent
Members of RAM group will meet every Friday from 4-6 pm.
On Tuesday, need to talk to Miss Polen. I need to start recording hours and get 10 to 15 time sheets to start recording hours.

\labday{Monday, July 1, 2019}
From the meeting with Dr. Miah:
This is what needs to happen before scheduling the meeting with Ashraf: 
\begin{enumerate}
\item Copy and paste the python script implementing the plugload api and control the WeMo switch
\item Document the current tkinter GUI used to control the motor
\end{enumerate}
In the meeting with Ashraf:
\begin{enumerate}
\item Previously mentioned items
\item Discuss ways to improve the current agent based system in BEMOSS
\end{enumerate}

\labday{Wednesday, July 3, 2019}
In a completely separate file, remove all the device api python scripts and try to implement the wemo plug in a new directory. Comment out all functions as needed. Also, the next meeting may be delayed until friday if Eric is not available as well as Caleb. 

\labday{Friday, July 5, 2019}
Deadline for DC motor integration has been changed to July 31; however, I would like to have everything done in a week.
\begin{itemize}
\item
Caleb's presentation:\\
\LaTeX is good for mathematics and presenting with page numbers
\item
Amr's presentation:\\
Had problems modeling joint\\
What is URDF file?\\
\end{itemize}

\labday{Friday, July 19, 2019}
For next pres:\\
Add slide for VOLTTRON (need to thouroughly understand it).\\
IPE figure for remote GPIO next to text.\\

Questions\\
Caleb's pres:\\
Is that a voltage regulator on Issues/Pictures slide?\\
It is a 7805 regulator.\\
Does the Beaglebone have GPIO?\\
Yes\\
What is the API mode?\\
Can easily configure modules.
\bigbreak\noindent
Amr's pres:\\
Can you explain what URDF is?\\
XML file which contains physical representation of the robot.

\labday{Tuesday, July 30, 2019}
In today's short meeting with Dr. Miah:
\begin{itemize}
\item Send Ashraf an email with the screenshot of the logging error in BEMOSS, see how he can resolve the problem
\begin{itemize}
\item We must submit work without any errors in the software
\end{itemize}
\item Work on the issue tonight
\item I should know the INs and OUTs of BEMOSS by now; however, I have failed to complete this task
\begin{itemize}
\item Without any hesitation, I should know where to implement pieces of software and not have to guess.
\end{itemize}
\item Dr. Miah will send me email with resources on software design
\item Deadlines are absolutely crucial and they cannot be missed
\item I need to learn object-oriented programming thoroughly
\end{itemize}

\labday{Friday, August 2, 2019}
From my pres:\\
Email Ashraf about issues. Stick with current version of BEMOSS.\\
Caleb's pres:\\
How does pc use powershell to communicate with Zigbee boards?\\

\labday{Monday, August 12, 2019}
Questions for meeting with Ashraf:
\begin{itemize}
\item How to add a new device type?
\item What other devices might be useful in adding?
\end{itemize}
In pgadminIII,
Look into\\
Supported Device list table\\
Not looking adding other types of devices

\labday{Monday, August 19, 2019}
Questions for meeting with Dr. Miah:
\begin{itemize}
\item When should I have the paper done by?
\item Should I draw a block diagram for the motor controller circuit in the lab?
\item Must the new device be a different type of device or be from a new category?
\end{itemize}

\labday{Friday, August 30, 2019}
Questions for Dr. Miah in meeting with 

\labday{Friday, September 13, 2019}
One possible idea is to implement in BEMOSS is default 'OFF' state.
In the group meeting, Eric suggested adding a way to set the default state of a device to the 'ON' state. This could be useful as some loads require initialization in the 'ON' state.

\labday{Tuesday, September 17, 2019}
Questions for Dr. Miah in the 11:00 meeting:
\begin{itemize}
\item What is the difference between a building energy management system and a building automation system?
Building energy management is utilizing policies and algorithms to help both meet energy demand in facilities and lower energy demand. On the other hand, building automation system is the centralized control of different building systems such as HVAC and lighting systems. Devices may be scheduled to turn on at certain times of the day or week. Data may be collected over time for analysis.
\item Should I be researching other open source/proprietary building automation systems or building energy management systems?
Right now, the platforms that were mentioned in the paper are fine. However, I do need to do a background study on some different platforms. They need to be briefly discussed in the problem description.

In the future, we are looking to build a mobile application rather than a web application for the purpose of building energy management. 
\end{itemize}
 
\labday{Tuesday, October 8, 2019}
Today, Dr. Miah shared the news that the university cancelled funding for all internal grant award programs, so the proposal will not be submitted to the university for funding. Instead, now the focus will be brought back to working on adding the Kill A Watt meter to BEMOSS and implementing the DC motor speed control algorithm. By November 15, research for the Kill A Watt meter needs to be completed (i.e. need to determine what parts to purchase so that Mr. Mattus can purchase the materials). The speed control algorithm will be started on Dec 1 and continue through February 2020. The gantt charts for 2019 and 2020 detailing the work are shown below.

\begin{figure}[H]
%\begin{sidewaysfigure}
    \begin{ganttchart}[hgrid,
    vgrid={*4{black, dotted},*3{white},*1{black,
    dotted},*3{white},*1{black, dotted},*3{white}},
    group/.append style={draw=black, fill=yellow!50},
    group left shift = 0,
    group right shift = 0,
    y unit title=.8cm,
    %y unit chart=.8cm,
    milestone label font=\tiny,
    bar label font=\small, 
    group label font=\small,
    bar/.append style={draw=black,fill=blue},
    bar left shift= 0,
    bar right shift= 0,
    bar height = .2]{1}{16}
    \gantttitle{2019}{16}\\
    \gantttitle{Sep}{4}
    \gantttitle{Oct}{4}
    \gantttitle{Nov}{4}
    \gantttitle{Dec}{4}\\
    % \ganttgroup{DC motor speed control algorithm}{5}{8}\\
    \ganttgroup{Research new device}{5}{12}\\
    \ganttbar{Research device}{5}{10}\\
    \ganttbar{Implemention in BEMOSS}{11}{12}\\
    \ganttgroup{DC motor speed control}{13}{15}\\
    \ganttbar{Research algorithm}{13}{15}
    \end{ganttchart}
\caption{Gantt chart for Fall 2019}
\label{gantt1}
%\end{sidewaysfigure}
\end{figure}
\begin{figure}[H]
%\begin{sidewaysfigure}
    \begin{ganttchart}[hgrid,
    vgrid={*4{black, dotted},*3{white},*1{black,
    dotted},*3{white},*1{black, dotted},*3{white},*1{black,dotted},*3{white}},
    group/.append style={draw=black, fill=yellow!50}, 
    group left shift = 0,
    group right shift = 0,    
    y unit title=.8cm,
    %y unit chart=.8cm,
    milestone label font=\tiny,
    bar label font=\small, 
    group label font=\small,
    bar/.append style={draw = black, fill=blue},
    bar left shift= 0,
    bar right shift= 0,
    bar height = 0.2,
    bar incomplete/.append style={fill=red}]{1}{20}
    \gantttitle{2020}{20}\\
    \gantttitle{Jan}{4}
    \gantttitle{Feb}{4}
    \gantttitle{Mar}{4}
    \gantttitle{Apr}{4}
    \gantttitle{May}{4}\\
    \ganttgroup{DC motor speed control (cont.)}{2}{8}\\
    \ganttbar{Implementation in BEMOSS}{2}{8}\\
    \ganttgroup{Machine learning algorithm}{9}{18}\\
    \end{ganttchart}
\caption{Gantt chart for Spring 2020}
\label{gantt1}
%\end{sidewaysfigure}
\end{figure}
After a second meeting with Dr. Miah where I presented to him the information I found about the Kill A Watt meter(s), we decided against implementing the Kill A Watt meter as the whole motivation for adding it to BEMOSS is to wirelessly measure the electricity demand (W or kW) of electrical/electronic loads. I did not mention to him the fact that the WeMo Insight switch can already wirelessely measure this quantity and is fully supported by BEMOSS. Thus the the Kill A Watt would be unecessary. Instead, the plan is to use the existing switch to measure power consumption of a new IoT device that is not currently supported by BEMOSS. The charts above have been updated with this information.
\end{addmargin}
%----------------------------------------------------------------------------------------
%	BIBLIOGRAPHY
%----------------------------------------------------------------------------------------


\bibliographystyle{plain}
\bibliography{bib/seniorProject2017.bib}


% \begin{thebibliography}{9}

% \bibitem{lamport94}
% Leslie Lamport,
% \emph{\LaTeX: A Document Preparation System}.
% Addison Wesley, Massachusetts,
% 2nd Edition,
% 1994.

% \end{thebibliography}

%--------------------------------------------------------------------------------------
\end{document}


%%% Local Variables:
%%% mode: latex
%%% TeX-master: t
%%% End:
